\begin{proposition}[E.39 --- Proved: the exact torsion lost by the compact and unit path comparisons]
\label{E-integral-path-comparison}
Retain the cubic lattice $V$, the cusp order and the matrix $U$
of (E18.8)--(E18.9), and put $K=\mathbb Q(\zeta_3)$.
For each nonorigin cusp $d$, let $f_d=L_d/L_O$ be the
normalized tangent-line unit of E.36, and let
$\ell_d:V\to\mathbb Z_3(1)$ be its geometric Kummer map.
Write $T=T_3E$ for the compact Tate module. Give
$B=\mathbb Z_3(1)^8$ the Galois action which permutes the
eight unit coordinates and multiplies by the cyclotomic
character. There is an actual equivariant injection
\[
 j:V\longrightarrow T\oplus B,\qquad
 v\longmapsto\bigl(p(v),(\ell_d(v))_d\bigr),\qquad
 j=\begin{pmatrix}I_2&0\\-U&3I_8\end{pmatrix}.
 \tag{E.39.1}
\]
Its cokernel $Q$ is $(\mathbb Z/3)^8$. The action on $Q$
over $K$ is trivial; over $\mathbb Q$, complex conjugation
acts by minus the permutation
\[
 (A_1\ A_2)(B_1\ B_2)(C_1\ C_2),
 \tag{E.39.2}
\]
with $A_0,B_0$ fixed by the permutation.
For continuous cohomology of the full absolute Galois
groups, and also of their quotients unramified away from
three and infinity, the exact kernels are
\[
 \begin{aligned}
 \ker H^1(G_K,j)&\simeq(\mathbb Z/3)^8,\\
 \ker H^1(G_{\mathbb Q},j)&\simeq(\mathbb Z/3)^3.
 \end{aligned}
 \tag{E.39.3}
\]
These isomorphisms are the connecting maps from $Q^{G_K}$
and $Q^{G_{\mathbb Q}}$, respectively. In particular, a
nonempty fibre of this cohomological comparison has
exactly $3^8$ elements over $K$ and $3^3$ over
$\mathbb Q$. Compact Kummer classes and all eight unit
cohomology classes therefore do not determine an integral
abelianized path cohomology class. One unresolved datum
in the integral version of E.36 is the actual class's
position in this explicitly specified fibre.
\end{proposition}
\begin{proof}
The divisor of $f_d$ is $3(d-O)$, so its restriction to
the boundary lattice is the row $3e_d^*$. The compact
rational splitting of E.18 has columns $\binom{3I_2}{U}$.
The Kummer map of a geometric unit vanishes on that
compact summand: equivariance at a good Frobenius compares
its weight-one compact eigenvalues with the weight-two
cyclotomic eigenvalue. Thus
\[
 \ell_d=(-U_d,3e_d^*),
\]
which proves (E.39.1) with its actual unit and compact
maps. Galois equivariance follows either from these maps
or directly from the matrices (E18.9). Subtracting the
lower left block using integral row operations gives
$\operatorname{diag}(I_2,3I_8)$. The quotient map can be
written
\[
 T\oplus B\longrightarrow Q,\qquad
 (h,b)\longmapsto b+Uh\pmod3.
 \tag{E.39.4}
\]
Over $K$, the cusp permutation is trivial and
$\chi_3\equiv1\pmod3$. Applying (E18.9) to (E.39.4)
therefore gives the trivial action. Over $\mathbb Q$,
the same matrix identity gives (E.39.2).

The invariant subgroup $(T\oplus B)^G$ is zero for each
of the four indicated arithmetic groups. For $T$ this
follows by choosing a good Frobenius whose eigenvalues
have complex absolute value greater than one. For $B$,
first pass to $K$, where it is a direct sum of
$\mathbb Z_3(1)$; the image of the cyclotomic character
is infinite, so no nonzero vector is fixed. These proofs
apply to the stated unramified quotients because the
actual representations factor through them.
The continuous cohomology sequence of
\[
 0\longrightarrow V\xrightarrow{j}T\oplus B
       \longrightarrow Q\longrightarrow0
\]
then identifies $Q^G$ injectively with the kernel in
degree one. The finite quotient admits a continuous
set-theoretic section, so this is the ordinary long exact
sequence for these continuous coefficient modules.
The fixed space of minus the permutation (E.39.2) has
the three basis vectors
\[
 e_{A_1}-e_{A_2},\qquad e_{B_1}-e_{B_2},\qquad
 e_{C_1}-e_{C_2}.
\]
The two fixed permutation coordinates contribute no
fixed vectors in characteristic three. This proves
(E.39.3) and the fibre cardinalities.

The connecting classes can also be read as explicit
cocycles. Let $t_d$ be the $d$th boundary basis vector
of $V$, and lift the $d$th basis vector of $Q$ rationally
by $t_d/3$. Its connecting cocycle is
\[
 b_d(g)=g(t_d/3)-t_d/3.
 \tag{E.39.5}
\]
Over $K$ this is integral and, on the actual free
arithmetic generators of E.38, has
$b_d(\tau)=b_d(\tau')=2t_d$. The three displayed
differences give the invariant classes over
$\mathbb Q$, with value zero on complex conjugation.
Their nonzero order three follows from the injective
connecting map; multiplying (E.39.5) by three gives
the coboundary of $t_d$.

The program \texttt{integral\_path\_comparison.py}
checks the exact coefficient maps, both arithmetic
equivariance identities, all ten Smith factors, the
three invariant vectors, and these cocycle values.
In the attempted integral automorphism calculation,
substituting arbitrary representatives of the compact
and unit cohomology classes can change the result by
(E.39.5). No such representative is asserted to be the
actual tangential path. The remaining comparison must
retain the correlated choices of compact division
points and unit roots. The manuscript anchor is
\S51.16, folios733v--736v, scans172--178, and
LM7-E.23(i).
\end{proof}
